\documentclass[tikz,border=8pt]{standalone}
\usepackage{ctex}
\usepackage{amsmath}
\usepackage{amssymb}
\usetikzlibrary{calc, arrows.meta, positioning}
\begin{document}
\begin{tikzpicture}[
  >={Stealth[length=2.6mm]},
  mode/.style={draw, thick, rounded corners=3pt, align=center, font=\small,
               text width=3.4cm, inner sep=4pt, minimum height=1.05cm},
  imp/.style={->, very thick, blue!60!black},
  tag/.style={draw, thin, rounded corners=2pt, align=left, font=\scriptsize,
              fill opacity=0.9, text opacity=1},
]
  % 标题
  \node[font=\bfseries] at (3.0,4.7)
    {随机收敛的四种模式与强弱层次（$\Rightarrow$ 表示「蕴含」，反向一般不成立）};

  % 四个收敛模式
  \node[mode, fill=blue!8]   (as) at (0.0,3.2)
    {\textbf{几乎必然 a.s.}（最强）\\[1pt] 每条样本路径都收敛};
  \node[mode, fill=cyan!8]   (l2) at (6.0,3.2)
    {\textbf{均方 / $L^p$ 收敛}\\[1pt] 平方偏差的期望 $\to 0$};
  \node[mode, fill=green!10] (p)  at (3.0,1.4)
    {\textbf{依概率 P}（居中）\\[1pt] 犯大错的概率 $\to 0$};
  \node[mode, fill=orange!12](d)  at (3.0,-0.6)
    {\textbf{依分布 d}（最弱）\\[1pt] 仅分布形状趋同 $F_n\to F$};

  % 蕴含箭头
  \draw[imp] (as) -- (p) node[midway, above left=-2pt, font=\footnotesize, blue!60!black] {$\Rightarrow$};
  \draw[imp] (l2) -- (p) node[midway, above right=-2pt, font=\footnotesize, blue!60!black] {$\Rightarrow$};
  \draw[imp] (p)  -- (d) node[midway, right=1pt, font=\footnotesize, blue!60!black] {$\Rightarrow$};

  % 极限定理落点标注
  \node[tag, fill=green!12] (lln) at (-2.0,0.4)
    {\textbf{大数定律 LLN}\\ $\bar X_n \to \mu$（常数）\\ 落在 P / a.s. 层};
  \draw[->, gray, thin] (lln.east) to[bend right=10] (p.west);

  \node[tag, fill=orange!14] (clt) at (7.7,-0.4)
    {\textbf{中心极限定理 CLT}\\ 标准化和 $\to \mathcal{N}(0,1)$\\ 落在 d 层};
  \draw[->, gray, thin] (clt.west) to[bend left=10] (d.east);

\end{tikzpicture}
\end{document}
